\chapter{Related Work}
\label{ch:related}

This chapter surveys the research landscape across five areas directly relevant
to this thesis: biometric template protection, fuzzy vault schemes, multimodal
biometric fusion, iris recognition and fuzzy commitment, and blockchain-based
biometric systems. A gap analysis table at the end positions this work against
the state of the art.

\section{Biometric Template Protection}

Biometric Template Protection (BTP) approaches fall into two broad categories:
cancelable biometrics and biometric cryptosystems.

\textbf{Cancelable biometrics} \cite{ratha2007generating} apply intentional,
repeatable but non-invertible transformations to raw biometric features before
storage. The transformed template is stored; the original features are never saved.
If the template is compromised, a different transformation produces a fresh,
unlinkable template from the same biometric. Teoh et al.\ \cite{teoh2004random}
introduced random multispace quantization for face and palmprint data. The primary
criticism of cancelable biometrics is that the transformation may reduce
discriminability.

\textbf{Biometric cryptosystems} \cite{jain2008biometric} bind a cryptographic
key to biometric data such that the key can only be reconstructed when a query
biometric is sufficiently similar to the enrolled one. Dodis et
al.\ \cite{dodis2008fuzzy} formalized this binding as \emph{fuzzy extractors},
providing information-theoretically secure constructions with error-tolerance.
The most widely studied instantiations are the fuzzy vault \cite{juels2002fuzzy}
and the fuzzy commitment \cite{hao2006combining}.

Rathgeb and Uhl \cite{rathgeb2011survey} provide a comprehensive survey of
both paradigms, concluding that biometric cryptosystems offer stronger provable
security guarantees but face greater challenges with noisy, unaligned data. This
thesis works within the biometric cryptosystem paradigm.

\section{Fuzzy Vault Schemes}

The fuzzy vault \cite{juels2002fuzzy} was proposed by Juels and Sudan (2002) as
a cryptographic construction for binding a secret to an unordered, fuzzy set.
A degree-$k$ secret polynomial $p(x)$ is evaluated at biometric-derived points;
these genuine points are hidden among chaff that do not lie on $p(x)$.
Decoding requires at least $k{+}1$ genuine matches for Lagrange interpolation.

Clancy et al.\ \cite{clancy2003secure} first applied the fuzzy vault to
fingerprint minutiae. Uludag et al.\ \cite{uludag2005fuzzy} added alignment
helper data, achieving EER 4.0\% on FVC2002~DB2. Nandakumar et
al.\ \cite{nandakumar2007fuzzy} improved alignment via high-curvature minutiae,
reducing EER to approximately 2\%.

Tams \cite{tams2013attacks} demonstrated record-multiplicity attacks: multiple
vaults from the same user can be cross-matched to identify genuine points. This
underscores the need for enrollment-independent structures that ensure vault
points are unlinkable across enrollments---a property provided by the ECDLP-based
construction of this thesis.

Maurya et al.\ \cite{maurya2025enhancing} substantially advanced the state of
the art by replacing GF($2^{16}$) encoding with elliptic curve cryptography over
Brainpool~P256r1 \cite{rfc5639}. Minutiae coordinates are quantized and mapped to
EC points via $P_i = q_i \cdot G$, a one-way mapping protected by ECDLP with
security strength $\sim\!2^{128}$ operations. Their scheme achieves EER as low as
0.005\% on FVC2002~DB1 with vault entropy exceeding 7.7~bits per coordinate. This
ECC vault is the unimodal baseline for the present work.

\section{Multimodal Biometric Fusion}

Multimodal biometrics improves recognition accuracy and spoof resistance by
combining evidence from multiple traits or sensors
\cite{ross2003information,jain2004multibiometrics}. Score-level fusion is most
common in unconstrained systems. Inside a fuzzy vault, however, fusion must be
integrated at the feature or vault-construction level because no matching score
is ever explicitly computed.

Nandakumar and Jain \cite{nandakumar2008multibiometric} fused fingerprint and
iris within a single fuzzy vault by including genuine points from both modalities.
Both modalities must successfully contribute to the polynomial reconstruction,
creating the AND-fusion constraint: $\mathrm{GAR}_\mathrm{AND} =
\mathrm{GAR}_\mathrm{FP} \times \mathrm{GAR}_\mathrm{Iris} \leq
\mathrm{GAR}_\mathrm{FP}$.

Moon et al.\ \cite{moon2010multibiometric} applied the same AND-fusion principle
to face and fingerprint, with the face module remaining mandatory. Kelkboom et
al.\ \cite{kelkboom2011multialgorithm} fused multiple fingerprint algorithms within
a fuzzy commitment framework but did not address cross-modal fusion with a
second biometric trait.

Three additional tight-coupling strategies evaluated in this thesis---polynomial
blinding (Architecture B), dense chaff seeding (Architecture C), and
AES-256-GCM vault encryption (Architecture D)---all require iris key recovery for
decoding and exhibit the same AND-fusion penalty. This thesis is the first to break
this structural pattern by treating the iris as a booster: contributing when
available, falling back to unimodal when not.

\section{Iris Recognition and Fuzzy Commitment}

Daugman \cite{daugman2004how} established the modern framework for iris
recognition. The iris is encoded as a binary iris code using multi-scale 2-D
Gabor wavelets, with fractional Hamming distance (FHD) as the matching metric.
FHD is near 0.10--0.15 for same-iris pairs and near 0.46 (random) for
different-iris pairs, providing wide separation. This makes iris one of the most
discriminative biometric traits when captured under near-infrared illumination.

Hao et al.\ \cite{hao2006combining} introduced the fuzzy commitment scheme for
iris codes, binding a 140-bit cryptographic key to a 2048-bit iris code using
error-correcting codes. Up to 25\% bit-level error can be corrected. This work
established that iris codes can be used reliably in fuzzy cryptosystems with
appropriate error-correcting codes.

Hollingsworth et al.\ \cite{hollingsworth2009fragile} identified fragile bits in
iris codes---positions near phase boundaries where encoding is unstable---and showed
that masking them significantly improves matching accuracy.

In this thesis, a repetition-code fuzzy commitment is applied to the full
131{,}072-bit iris code from CASIA-Iris-Interval images. Block-wise majority
decoding with bit interleaving (distributing spatially correlated errors across
blocks) achieves a key recovery GAR of approximately 62\% at impostor FAR of
0.10\%, with block size $B_s = 255$~bits.

\section{Blockchain-Based Biometric Systems}
\label{sec:blockchain_related}

The intersection of biometrics and blockchain has attracted growing interest as a
potential solution to the centralized storage problem.

Delgado-Mohatar et al.\ \cite{delgado2020blockchain} surveyed the opportunities
and challenges of storing biometric templates on blockchain, identifying high
on-chain storage costs, privacy risks from immutable records, and latency as the
key barriers. Their key recommendation---off-chain storage with on-chain hash
anchoring---is the architecture adopted by D-IBFV.

Acquah et al.\ \cite{acquah2020securing} confirmed that off-chain storage with
on-chain hash anchoring is the scalable pattern and that on-chain template storage
incurs prohibitive gas costs at scale.

Goel et al.\ \cite{goel2019securing} secured CNN-based biometric feature extraction
via blockchain with Shamir-based key recovery. However, unprotected feature data
is stored on-chain, providing process auditability without template-level
cryptographic protection.

Lee and Jeong \cite{lee2021securing} proposed BDAS, a blockchain-based
decentralized authentication system that fragments templates across a distributed
ledger. BDAS achieves decentralization but without a fuzzy vault construction or
formal protection guarantees.

Sharma et al.\ \cite{sharma2024multimodal} combined fingerprint and iris via
cancelable transforms on blockchain (EER~$\sim$8\%). Sharma and Dwivedi
\cite{sharma2026fabric} deployed a similar approach on Hyperledger Fabric, reporting
sub-second latency, without vault-level splitting.

Yun et al.\ \cite{yun2024biorollup} proposed Bio-Rollup, integrating zk-SNARK
auditing with layer-2 blockchain scaling for biometric query privacy, but without
cryptographic vault protection for stored templates.

Prakasha et al.\ \cite{prakasha2025privacy} surveyed privacy-preserving biometric
authentication on blockchain, highlighting the need for stronger template protection,
scalable architectures, and formal security proofs.

\section{Summary and Gap Analysis}

Table~\ref{tab:gap} summarizes the key characteristics of related approaches.
Three gaps are directly addressed by this thesis:

\begin{enumerate}
  \item \textbf{No prior multimodal fuzzy vault is penalty-free.} All existing
  multimodal schemes gate decoding on the auxiliary modality.

  \item \textbf{No prior blockchain biometric system applies vault-level secret
  sharing with formal proofs.} Existing systems store templates as monolithic
  objects.

  \item \textbf{No prior system demonstrates zero accuracy degradation from
  decentralization.} D-IBFV proves and empirically confirms bit-exact reconstruction.
\end{enumerate}

\begin{table}[h]
\centering
\caption{Gap analysis: multimodal and blockchain-based biometric schemes.}
\label{tab:gap}
\small
\begin{tabular}{@{}lcccc@{}}
\toprule
\textbf{Scheme} & \textbf{Aux.\ Required} & \textbf{Zero-Penalty} &
\textbf{Decentralized} & \textbf{Formal Proofs} \\
\midrule
Nandakumar \cite{nandakumar2008multibiometric} & Yes & No & No & No \\
Moon \cite{moon2010multibiometric}             & Yes & No & No & No \\
Kelkboom \cite{kelkboom2011multialgorithm}     & Yes & No & No & No \\
Arch.\ A--D (this work, compared)             & Yes & No & No & No \\
\midrule
Delgado-Mohatar \cite{delgado2020blockchain}   & --- & --- & Yes & No \\
Sharma \cite{sharma2024multimodal}             & Yes & No & Yes & No \\
Lee BDAS \cite{lee2021securing}                & --- & --- & Yes & No \\
Yun Bio-Rollup \cite{yun2024biorollup}         & --- & --- & Yes & No \\
\midrule
\textbf{IBFV (this thesis)}   & \textbf{No} & \textbf{Yes} & No & \textbf{Yes} \\
\textbf{D-IBFV (this thesis)} & \textbf{No} & \textbf{Yes} & \textbf{Yes} & \textbf{Yes} \\
\bottomrule
\end{tabular}
\end{table}
